\section{相关工作}

本文落在四支文献的交汇处。每支既给方法论提示，也给出需要绕开的坑。

\textbf{网络度量的一支。} 中心性族与秩相关口径由 \citet{freeman1978} 定形，
\citep{wasserman1994} 把它整理成可操作的教科书体系；本篇的度数、出场章节数与
Spearman $\rho$ 都按这一支的定义算，并列秩取平均秩。要绕开的坑是：中心性在稀疏图上不稳，
而本篇全图有 537 个孤立点、mean degree 只有 6.27，任何单点驱动的结论都要降级为描述统计量——
这正是第 6 节 L4 与 L6 的由头。

\textbf{历史学用正式网络方法的一支。} \citet{lemercier2015} 论证历史学者为什么、
以及怎样用网络方法，是本篇的方法论许可。她同时点明了最要紧的坑：从历史文本抽出来的「关系」
是叙述者建构出来的关系，不是被研究时代里发生过的关系。本篇全部结论因此只落在
「这部叙事把谁和谁写成什么关系」，一次也不外推到明代社会结构。

\textbf{解释关系之功能的一支，本篇借它而不用它。} \citet{granovetter1973} 的弱关系、
\citet{burt1992} 的结构洞、\citet{coleman1988} 的社会资本，讲的是义务、闭合与投资：
礼物怎么生产关系，宗族内部为什么高度闭合。这些内容给了本篇解释框架——「名分」在书中
是否正是这种闭合的技术。但要绕开的坑是理论比证据跑得远：本篇量的是
「关系被写成了哪个类别」，不是「关系起了什么功能」，因此不给任何机制断言。

\textbf{通俗叙事与词汇统计的一支。} \citet{michel2011} 用海量数字化图书做数量文化史，
证明了词频可以成为历史证据，也顺手证明了它的边界：词频不等于语义场，
一个词的高频经常只是行文习惯。\citet{brook1998} 研究明代的物质文化与礼物流动，
给了「馈赠—庇护」这条线的历史纵深。本篇的坑是别把词频当成网络：Michel 那一支停在词这一层，
本篇往下走一层，问词所指的人在图上是怎样的形状。

\textbf{本文的位置。} 四支里，各有其用：或给度量，或给对照假设，或给规模，或给对象，但没有一支把「某一组这个词」与「关系类别构成」接起来：现有研究要么讨论该领域的演变，要么统计全书词频，中间缺一层——某组词汇在网络上对应什么形状。本文用主题词命中加 egocentric 子图类别构成补上这一层，每个数字都能回到可复算的底稿。
